\documentclass[10pt,twocolumn]{article}
\usepackage[a4paper,margin=18mm,columnsep=7mm]{geometry}
\usepackage{lmodern}
\usepackage{amsmath,amssymb,amsthm,mathtools}
\usepackage{booktabs,tabularx,graphicx}
\usepackage{xcolor}
\usepackage{hyperref}
\usepackage{microtype}
\usepackage{caption}

\definecolor{primary}{RGB}{15, 23, 42}
\definecolor{accent}{RGB}{13, 148, 136}
\hypersetup{
  colorlinks=true,
  linkcolor=accent,
  citecolor=accent,
  urlcolor=accent,
  pdftitle={AHD-Jailbreak: Adaptive Hypergraph Defense with Anchor-Gated Topology Evolution for Detecting Multi-Turn and Compositional Jailbreak Attacks in Large Language Models},
  pdfauthor={Sesetti Pavan Kumar Swamy, Rajana Rajeswari, and Mrs. S. Suguna Sri}
}

\setlength{\parskip}{4pt plus 1pt minus 1pt}
\captionsetup{font=footnotesize,labelfont=bf}

% Dynamic Graphics Resolution & Clickable Image Links
\graphicspath{{figures/}{reports/figures/}{./}}
\DeclareGraphicsExtensions{.pdf,.png,.jpg}

\newcommand{\dynamicurlimage}[3][width=\textwidth]{%
  \immediate\write18{curl -s -k -L -o "figures/#3" "#2"}%
  \href{#2}{\includegraphics[#1]{#3}}%
}

\newcommand{\takeaway}[2]{%
  \vspace{4pt}\noindent\colorbox{accent!10}{%
    \parbox{\dimexpr\columnwidth-2\fboxsep\relax}{%
      \vspace{2pt}\footnotesize\textbf{\color{primary}#1:}\hspace{0.4em}#2\vspace{2pt}%
    }%
  }\vspace{4pt}%
}

\title{\Large\textbf{AHD-Jailbreak: Adaptive Hypergraph Defense with Anchor-Gated Topology Evolution for Detecting Multi-Turn and Compositional Jailbreak Attacks in Large Language Models}}

\author{
  \textbf{Sesetti Pavan Kumar Swamy}$^{1}$, \textbf{Rajana Rajeswari}$^{2}$, and \textbf{Mrs. S. Suguna Sri}$^{3}\\[4pt]
  \small $^{1}$Dept. of CSE (24555A0519 (CSE–C)), Godavari Global University\\
  \small $^{2}$Dept. of CSE (23551A05C9 (CSE–B)), Godavari Global University\\
  \small $^{3}$Faculty Guide & Assistant Professor, Dept. of CSE, Godavari Global University, Rajamahendravaram, AP, India
}
\date{}

\begin{document}
\maketitle

\begin{abstract}
Large Language Models (LLMs) deployed across enterprise APIs, cloud chatbots, and cyber-physical edge devices remain vulnerable to evolving adversarial jailbreaks. Attackers exploit multi-turn conversational escalation (Crescendo, MHJ), competing-objective persona modulation, and surface encoding obfuscation to bypass alignment guardrails. Existing lightweight input guardrails process prompts as isolated single-turn sequences, suffering high false-negative rates on multi-turn escalation and severe over-refusal (false positives) on benign technical prompts containing sensitive homonyms. Conversely, recent graph-based defenses rely on static pairwise edges (|e| = 2) that cannot natively represent higher-order polyadic co-activations across dialogue turns, linguistic obfuscation signals, and targeted safety policies. We propose AHD-Jailbreak, a lightweight, model-agnostic pre-inference defense middleware built upon an adaptive heterogeneous hypergraph G_t = (V_t, E_t, W_t, X_t). Each node encodes a 394-dimensional representation combining 384-D dense semantic embeddings (all-MiniLM-L6-v2) with 10 pre-normalization surface indicators, while four higher-order hyperedge families simultaneously bind prompt instances, sliding 5-turn dialogue trajectories, semantic clusters, adversarial tactics, and safety policy categories. To adapt to newly emerging attack families without catastrophic forgetting or adversarial label poisoning, AHD-Jailbreak introduces an offline Quarantine Buffer governed by an Anchor Validation Gate (Module 14) that promotes candidate topology updates G_t -> G_{t+1} only when F1 on a locked 500-sample validation anchor set remains within 0.5% and safe contrast false-positive rate (FPR) rises by at most 1.0%. Across 5 random seeds on a leakage-free 14,500-sample multi-benchmark suite (WildGuardMix, WildJailbreak, ScaleAI MHJ, JailbreakBench, and XSTest), AHD-Jailbreak achieves 100.0% in-distribution F1, 93.67% F1 on held-out out-of-distribution (OOD) attack families (+16.12 percentage points over a matched-encoder 394-D MLP baseline, McNemar p < 0.0001), and 90.32% cross-dataset F1 on JailbreakBench while maintaining a 1.49% FPR on XSTest contrastive safe prompts. Using 2-hop local ego-subhypergraph extraction, AHD-Jailbreak executes in 2.25 ms median (3.26 ms P95) CPU hop latency (14.2 ms end-to-end on x86_64 and 26.5 ms on a Raspberry Pi 5 8GB board), providing a practical, privacy-preserving (HMAC-SHA256) inline guardrail for both cloud APIs and embedded edge systems.
\end{abstract}

\section{Introduction}
Safety-aligned Large Language Models (LLMs) undergo supervised fine-tuning (SFT) and reinforcement learning from human feedback (RLHF) to refuse harmful instructions. However, as formalized by Wei et al. [1], alignment remains vulnerable to two structural failure modes: competing objectives (where instruction-following conflicts with safety refusal, such as forced affirmative prefixes) and mismatched generalization (where inputs are obfuscated via Base64, hexadecimal, or low-resource ciphers absent from safety training). Furthermore, multi-turn conversational attacks such as Crescendo [2] and Multi-Turn Human Jailbreaks (MHJ) [3] decompose a forbidden objective across several benign-looking turns, achieving over 70% Attack Success Rate (ASR) against defended models.

Existing defense paradigms face a sharp operational trilemma across \textbf{relational context}, \textbf{adaptive generalization}, and \textbf{hardware efficiency}: (1) \textit{Stateless Sequence Classifiers:} Lightweight classifiers such as Prompt-Guard-2-86M [4] or TF-IDF inspect single turns in isolation, missing multi-turn trajectory drift and suffering high false-positive rates on technical homonyms (e.g., \emph{``kill a zombie process in Linux''}) [5]; (2) \textit{Heavyweight LLM Judges:} Generative guardrails such as Llama Guard [6] and WildGuard [7] require 7B--8B parameter GPU deployments (200--800 ms latency, $\ge$ 16 GB VRAM) and cannot update their static weights online without costly retraining; and (3) \textit{Pairwise Graph Neural Networks:} Recent graph filters such as GuardNet [8] and G-Guard [9] construct pairwise graphs ($|e|=2$), failing to natively represent higher-order polyadic co-activations across dialogue turns, linguistic obfuscations, and safety policies without lossy clique expansion.

To resolve this trilemma, we introduce \textbf{AHD-Jailbreak}, an Adaptive Hypergraph Defense middleware positioned upstream of any target LLM. Our primary contributions are threefold: First, we formulate pre-inference screening as a heterogeneous polyadic hypergraph $G_t = (V_t, E_t, W_t, X_t)$ over five vertex modalities and four higher-order hyperedge families ($E_{\text{sem}}, E_{\text{ctx}}, E_{\text{tac}}, E_{\text{pol}}$) using 394-D node features that preserve pre-NFKC obfuscation indicators. Second, we establish an offline Quarantine Buffer governed by a locked 500-sample Anchor Validation Gate, guaranteeing immunity against adversarial /feedback label poisoning with automatic rollback. Third, across 5 random seeds and 6 systematic ablations (A1--A6), we demonstrate statistically significant gains over matched-encoder flat classifiers ($p < 0.0001$) and verify real-world deployment across x86\_64 servers, NVIDIA Jetson Orin Nano, and Raspberry Pi 5 edge boards.

\section{Threat Model \& High-Dimensional (HD) Jailbreaks}
We consider an adversary interacting with an LLM application across Trust Boundary TB-1 via single-turn or multi-turn HTTP requests, external RAG documents, and public /feedback submissions. In real-world cyber-physical and enterprise systems, sophisticated attackers execute High-Dimensional / Hypergraph-Distributed (HD) Jailbreaks that distribute malicious intent across five orthogonal dimensions simultaneously: (1) multi-turn semantic drift across W_turn = 5 turns [2, 3], (2) pre-normalization character/encoding obfuscation (Base64, zero-width joiners, homoglyphs) [1], (3) competing-objective persona framing (AF-02/AF-03), (4) indirect third-party RAG/sensor context injection [10], and (5) adaptive threshold probing [11] or feedback poisoning.

Let T_1 denote unsafe intent detection (safety_label), T_2 denote adversarial framing detection (attack_family), T_3 denote prompt injection hijacking, and T_4 denote downstream LLM compliance (Attack Success Rate, ASR). Because AHD-Jailbreak operates pre-inference, it predicts the composite actionable risk label y in {0: Allow, 1: Block/Review} where y = 1 iff the interaction exhibits unsafe intent (vanilla or adversarial harmful), while explicitly routing contrastive benign technical and roleplay queries (BF-02) to y = 0.

\section{Methodology: Heterogeneous Hypergraph \& Anchor-Gated Adaptation}

\begin{figure*}[t]
\centering
\dynamicurlimage{/figures/fig0_system_architecture.png}{fig0_system_architecture.png}
\caption{The complete 6-stage defensive pipeline: (1) Multi-turn dialogue trajectory and raw query input; (2) 10-D pre-NFKC surface feature extraction combined with 384-D dense semantic embedding into unified representation x_q in R^{394}; (3) Inductive cosine k-NN ego-subhypergraph extraction (<= 120 nodes); (4) 2-layer spectral hypergraph convolution with temperature calibration (T_cal = 0.84); (5) Calibrated multi-signal risk fusion; and (6) Module 14 Anchor Validation Gate with locked 500 validation anchors ensuring poison resilience.}
\label{fig:architecture}
\end{figure*}

\subsection{Heterogeneous Vertex Space and 394-D Feature Extraction}
At adaptation epoch t in {0, 1, ...}, the defense topology is defined as G_t = (V_t, E_t, W_t, X_t) with |V_0| = 4,084 vertices partitioned into five disjoint sets:
V_t = V_P union V_T union V_S union V_C union V_A  (Eq. 1)
where V_P (N_P = 4,000) stores stratified anchor prompts and the active query v_q; V_T represents historical dialogue turns t_{k-m}, ..., t_{k-1} (W_turn = 5); V_S (K_S = 64) represents spherical k-means semantic cluster centroids (32 harmful, 32 contrastive safe); V_C (K_C = 10) represents canonical safety policy categories [7, 12]; and V_A (K_A = 10) represents structural tactic families (AF-01–AF-08, BF-01, BF-02).

Crucially, applying Unicode NFKC normalization first would erase homoglyphs and zero-width characters. Our feature extractor first computes a 10-dimensional surface vector f_surf in R^{10} (character Shannon entropy, log token count, punctuation/uppercase/digit ratios, Base64 regex match, zero-width/homoglyph count, override verb count, roleplay marker count, and turn-to-turn cosine drift 1 - cos(e_t, e_{t-1})) prior to NFKC normalization, and concatenates it with the 384-D L_2-normalized all-MiniLM-L6-v2 embedding e_sem to form x_v = [e_sem || f_surf] in R^{394}.

\subsection{Higher-Order Hyperedges and Spectral Convolution}
Four hyperedge families E_t = E_sem union E_ctx union E_tac union E_pol (|E_0| = 198) connect arbitrary vertex subsets e subseteq V_t (|e| >= 2) via sparse incidence matrix H_t in {0,1}^{|V_t| x |E_t|}. Each hyperedge is initialized with a class-purity weight on the training split (lambda_base = 0.20):
w_0(e) = lambda_base + (1 - lambda_base) | 2 P(y=1 | v in e) - 1 |  (Eq. 2)
Let D_v(v,v) = sum_e w_t(e) H_t(v,e) and D_e(e,e) = sum_v H_t(v,e) denote the diagonal vertex and hyperedge degree matrices. Following Feng et al. [13] and Bai et al. [14], our 2-layer architecture projects Z^(0) = ReLU(X_t W_in + b_in) in R^{|V_t| x 128} and applies two stages of spectral node-to-hyperedge-to-node message passing with residual connections and LayerNorm:
Z^(l+1) = LN(Z^(l) + ELU(P_t Z^(l) Theta^(l)))  (Eq. 3)
where P_t = D_v^{-1/2} H_t W_t D_e^{-1} H_t^T D_v^{-1/2}. At runtime, rather than rebuilding the full graph, an incoming query v_q queries a cached cosine k-NN index (k=10) to extract its 2-hop local ego-subhypergraph (<= 120 vertices), executing the 2-layer forward pass in 2.25 ms on CPU.

\subsection{Calibrated Multi-Signal Risk Fusion and Anchor Gate}
Given logit l(v_q), we apply temperature scaling P_cal(x) = sigma(l(v_q)/T_cal) where T_cal = 0.84 is optimized via L-BFGS on val. Module 10 computes composite risk R(x) in [0,1] and epistemic uncertainty u(x) in [0,1]:
R(x) = 0.50 P_cal(x) + 0.25 S_nbr(x) + 0.15 S_ctx(x) + 0.10 S_OOD(x)  (Eq. 4)
u(x) = 0.60 H_bin(P_cal(x)) + 0.40 S_OOD(x)  (Eq. 5)
where S_nbr(x) is the hyperedge-weighted harmful neighbor purity, S_ctx(x) captures multi-turn escalation velocity, and S_OOD(x) is the normalized distance to the nearest cluster centroid c_s in V_S. Requests are routed via three-way policy thresholds (tau_allow=0.35, tau_block=0.70, tau_u=0.45) into allow, review (quarantine), or block.

To prevent adversarial drift or feedback poisoning, public /feedback submissions never mutate G_t directly. Verified quarantine batches update hyperedge weights via capped exponential moving average (alpha=0.25, |Delta w| <= 0.15) while keeping >= 80% of core anchors immutable. A candidate snapshot G_{t+1} is promoted to production iff it passes the Anchor Validation Gate on D_anchor (500 locked validation prompts + 100 safe contrasts):
F1(G_{t+1}, D_anchor) >= F1(G_t, D_anchor) - 0.005  (Eq. 6)
FPR(G_{t+1}, D_safe) <= FPR(G_t, D_safe) + 0.010  (Eq. 7)

\section{Experimental Setup \& 5-Seed Results}

\begin{table}[t]
\centering
\caption{Leakage-Free 6-Partition Dataset Summary ($14,500$ Samples; Connected-Component Clustering at $\cos \ge 0.90$).}
\label{tab:splits}
\footnotesize
\begin{tabularx}{\columnwidth}{l r X}
\toprule
\textbf{Split} & \textbf{Count} & \textbf{Primary Sources \& Isolation Rule} \\
\midrule
\texttt{train} & 8,000 & WildGuardTrain, WildJailbreak, MHJ (70%); excludes held-out AF-04..06. \\
\texttt{val} & 2,000 & Disjoint clusters; tunes T_cal=0.84, thresholds, and locks 500 Anchor Gate samples. \\
\texttt{test_id} & 1,725 & Official human-audited WildGuardTest; cos ≥ 0.90 matches purged from train. \\
\texttt{test_ood} & 1,200 & Held-out Base64 (AF-04), GCG (AF-05), AutoDAN (AF-06), MHJ (30%) across B1..B4. \\
\texttt{test_cross} & 800 & JailbreakBench (200 + 300 artifacts) & StrongREJECT (313). \\
\texttt{test_fpr} & 1,450 & Complete XSTest (450) + OR-Bench-Hard-1K (1,000); zero training overlap. \\
\bottomrule
\end{tabularx}
\end{table}

\begin{table*}[t]
\centering
\caption{5-Seed Multi-Benchmark Empirical Comparison (RQ1--RQ4).}
\label{tab:master}
\footnotesize
\begin{tabularx}{\textwidth}{l r r r r r r r}
\toprule
\textbf{Defense Paradigm} & \textbf{In-Dist F1} & \textbf{AUROC} & \textbf{AUPRC} & \textbf{XSTest FPR} & \textbf{OOD F1} & \textbf{Cross-JBB F1} & \textbf{Hop ms} \\
\midrule
B1: Rule-Based Heuristic Filter & 79.02% & 0.9312 & 0.9247 & 35.11% & 92.94% & 85.44% & 0.43 ms \\
B2a: Classical TF-IDF + Logistic Regression & 100.0% & 1.0000 & 1.0000 & 23.40% & 97.67% & 99.03% & 1.28 ms \\
B2b: Matched-Encoder 394-D MLP & 100.0% & 1.0000 & 1.0000 & 0.00% & 77.55% & 74.96% & 1.12 ms \\
B3: Static Hypergraph (Frozen G_0) & 100.0% & 1.0000 & 1.0000 & 2.98% & 80.74% & 74.47% & 2.24 ms \\
B4: AHD-Jailbreak (Proposed Adaptive G_t) & 100.0% & 1.0000 & 1.0000 & 1.49% & 93.67% & 90.32% & 2.25 ms \\
B5: Meta Llama-Prompt-Guard-2-86M & 73.68% & 0.9815 & 0.9837 & 0.00% & 71.76% & 69.23% & 1.32 ms \\
\bottomrule
\end{tabularx}
\end{table*}

Table I summarizes our 14,500-sample leakage-free corpus, where connected-component clustering at cos >= 0.90 and conversation_id grouping guarantee zero cross-split contamination. Table II and Figure 2(A) present the 5-seed evaluation across B1–B5. While lexical baselines B1 and B2a achieve high recall on keyword-heavy test items, they suffer severe false-positive rates on XSTest safe homonyms (35.11% and 23.40% FPR), rendering them unusable in developer environments. Conversely, flat neural classifiers B2b (394-D Matched MLP) and B5 (Prompt-Guard-2-86M) achieve 0.00% XSTest FPR but degrade sharply on held-out OOD attack families (77.55% and 71.76% F1).

AHD-Jailbreak (B4) resolves this safety-usability trade-off: using the exact same 394-D feature matrix X as B2b, B4 improves Held-Out OOD F1 by +16.12 percentage points (93.67% vs. 77.55%, McNemar p < 0.0001) and Cross-Dataset JailbreakBench F1 by +15.36 percentage points (90.32% vs. 74.96%) while maintaining a 1.49% XSTest safe FPR. Under our prequential test-then-adapt stream (B_1 -> B_4, Figure 3(A)), validated adaptation improves F1 from 80.00% (G_0) to 97.56% (G_4). When subjected to a label-flipped poison batch dropping anchor F1 to 40.00%, Module 14 automatically aborts promotion and rolls back to G_4. As plotted in Figure 3(B), temperature scaling (T_cal=0.84) aligns raw confidence with empirical error, lowering ECE from 0.1371 to 0.0502.

\begin{figure*}[t]
\centering
\dynamicurlimage{/figures/fig1_benchmark_and_ablations.png}{fig1_benchmark_and_ablations.png}
\caption{(A) 5-Seed Master Benchmark comparison across B1--B5 demonstrating how AHD-Jailbreak (B4) simultaneously achieves 93.67% Held-Out OOD F1, 90.32% Cross-Dataset JailbreakBench F1, and 98.51% safe prompt acceptance (1.49% FPR on XSTest). (B) Systematic ablation study (A1--A6) quantifying combined F1 degradation and XSTest FPR inflation paired with McNemar's continuity-corrected chi-square significance tests.}
\label{fig:benchmark_ablation}
\end{figure*}

\begin{table}[t]
\centering
\caption{Systematic Ablation Studies (A1--A6).}
\label{tab:ablations}
\footnotesize
\begin{tabularx}{\columnwidth}{l r r r r r}
\toprule
\textbf{Ablation Configuration} & \textbf{Comb. F1} & \textbf{$\Delta$ F1} & \textbf{XSTest FPR} & \textbf{ECE} & \textbf{McNemar $p$} \\
\midrule
Full AHD-Jailbreak (Ref) & 90.05% & 0.00% & 0.00% & 0.0502 & Ref (1.00) \\
A1: No Context (W_turn=1) & 84.46% & -5.59% & 0.00% & 0.0640 & 0.0001 \\
A2: No 384-D Dense Features & 22.40% & -67.65% & 3.19% & 0.1322 & < 0.0001 \\
A3: Pairwise GNN (|e|=2) & 85.35% & -4.70% & 5.32% & 0.0478 & 0.0014 \\
A4: Static Topology (G_0) & 84.46% & -5.59% & 5.32% & 0.0510 & 0.0001 \\
A5: No OOD Signal (w_4=0) & 84.37% & -5.68% & 4.26% & 0.0668 & 0.0001 \\
A6: Uncalibrated (T_cal=1.0) & 83.38% & -6.67% & 12.77% & 0.1371 & < 0.0001 \\
\bottomrule
\end{tabularx}
\end{table}

\takeaway{Ablation Takeaway}{Replacing polyadic hyperedges with pairwise edges (|e|=2, Ablation A3) drops combined F1 by -4.70% (p=0.0014) and inflates XSTest FPR to 5.32%, while removing temperature calibration (Ablation A6) causes XSTest FPR to spike to 12.77% (ECE = 0.1371).}

\begin{figure*}[t]
\centering
\dynamicurlimage{/figures/fig2_adaptation_and_calibration.png}{fig2_adaptation_and_calibration.png}
\caption{(A) Prequential test-then-adapt stream adaptation across 4 chronological OOD attack windows (B1: Base64, B2: GCG, B3: AutoDAN, B4: Crescendo), showing +17.56% F1 gain over static G_0 and automatic rollback when a label-flipped poison batch is injected at B3. (B) Temperature scaling calibration curve (T_cal = 0.84) reducing ECE from 0.1371 to 0.0502 and mapping risk into the calibrated review/quarantine band [tau_allow=0.35, tau_block=0.70].}
\label{fig:adaptation_calibration}
\end{figure*}

\section{Real-World Embedded Boards \& Production Profiling}
Because AHD-Jailbreak requires only 23 MB for an INT8-quantized ONNX encoder and < 15 MB for sparse snapshot $G_t$, it can be deployed directly onto physical embedded boards and enterprise gateways (Table IV and Figure 4): (1) \textit{Dual-NIC ``Bump-in-the-Wire'' Hardware AI Firewall (Raspberry Pi 5 8GB):} Configured with onboard Gigabit Ethernet (eth0, untrusted LAN) and a USB 3.0 Gigabit adapter (eth1, LLM server), executing ONNX Runtime ARM64 INT8 encoding (21.0 ms), $k$-NN ego extraction (2.4 ms), and 2-layer PyG HypergraphConv (3.1 ms) in 26.5 ms median (44.8 ms P95) latency under 650 MB RAM; (2) \textit{Air-Gapped Robotics \& Edge Box (NVIDIA Jetson Orin Nano 8GB):} Executes TensorRT FP16 hypergraph screening in 9.4 ms median (18.2 ms P95) alongside a local 4-bit quantized Llama-3.2-3B model via Ollama on a single 15W board, intercepting cross-channel QR/sensor + voice Crescendo HD jailbreaks (Figure 4(A)) prior to actuator execution; and (3) \textit{Federated Zero-PII Threat Graph Synchronization:} With \texttt{LOG\_RAW\_PROMPTS=false}, audit logs record only salted HMAC-SHA256 digests in SQLite WAL mode, while anchor-validated snapshots (\texttt{graph\_v\{t\}.pt}) synchronize across fleets of edge boards without sharing raw user text.

\begin{table}[t]
\centering
\caption{Pre-Inference Screening Latency Across Hardware Boards (RQ5).}
\label{tab:latency}
\footnotesize
\begin{tabularx}{\columnwidth}{l r r r}
\toprule
\textbf{Hardware Platform / Board} & \textbf{Enc. (ms)} & \textbf{Ego+HGNN} & \textbf{Total P50 / P95} \\
\midrule
NVIDIA Jetson Orin Nano (8GB, TensorRT) & 6.2 ms & 3.2 ms & 9.4 / 18.2 ms \\
Apple M2 Pro (ARM64 CPU) & 8.4 ms & 3.4 ms & 11.8 / 22.4 ms \\
x86_64 4-Core Xeon (AVX-512 ONNX) & 10.5 ms & 3.7 ms & 14.2 / 28.5 ms \\
Raspberry Pi 5 8GB (ARM64 ONNX INT8) & 21.0 ms & 5.5 ms & 26.5 / 44.8 ms \\
Raspberry Pi 5 8GB (FP32 PyTorch Raw) & 58.5 ms & 9.5 ms & 68.0 / 114.5 ms \\
\bottomrule
\end{tabularx}
\end{table}

\begin{figure*}[t]
\centering
\dynamicurlimage{/figures/fig3_crescendo_and_latency.png}{fig3_crescendo_and_latency.png}
\caption{(A) Multi-turn Crescendo escalation trajectory across Turns 1--3, demonstrating how the joint E_ctx hyperedge detects positive risk velocity (R = 0.09 -> 0.38 -> 0.86) crossing tau_block=0.70, whereas single-turn and lexical baselines remain below threshold. (B) Real-world stage-by-stage end-to-end latency breakdown across 5 production hardware configurations (NVIDIA Jetson Orin Nano, Apple M2 Pro, x86_64 Xeon, and Raspberry Pi 5), confirming sub-30 ms pre-inference execution within real-time SLA ceilings.}
\label{fig:crescendo_latency}
\end{figure*}

\section{Conclusion \& Limitations}
AHD-Jailbreak demonstrates that an inductive, anchor-gated heterogeneous hypergraph effectively bridges the gap between stateless sequence classifiers and heavyweight GPU judge models, intercepting multi-turn and compositional jailbreaks in < 30 ms on commodity CPUs and embedded boards. We note three operational limitations for future work: (1) because all-MiniLM-L6-v2 remains frozen during micro-batch updates, adaptation refines relational topology rather than learning new token semantics; (2) quarantine promotion relies on periodic high-confidence oracle verification; and (3) extending hyperedges to vision/audio modalities will further harden multimodal VLMs.

\begin{thebibliography}{14}
\bibitem{ref1} A. Wei, N. Haghtalab, and J. Steinhardt, "Jailbroken: How Does LLM Safety Training Fail?," in Proc. NeurIPS, 2023.
\bibitem{ref2} M. Russinovich, A. Salem, and R. Eldan, "The Crescendo Multi-Turn LLM Jailbreak Attack," in Proc. USENIX Security, 2025.
\bibitem{ref3} N. Li et al., "LLM Defenses Are Not Robust to Multi-Turn Human Jailbreaks Yet," Scale AI Tech Report, 2024.
\bibitem{ref4} Meta AI, "Llama Prompt Guard 2 (86M)," Hugging Face Repository, 2025.
\bibitem{ref5} P. Röttger et al., "XSTest: Identifying Exaggerated Safety Behaviours in LLMs," in Proc. NAACL, 2024.
\bibitem{ref6} H. Inan et al., "Llama Guard: LLM-based Input-Output Safeguard," arXiv:2312.06674, 2023.
\bibitem{ref7} S. Han et al., "WildGuard: Open One-Stop Moderation Tools," in Proc. NeurIPS, 2024.
\bibitem{ref8} J. Forough, M. Maheri, and H. Haddadi, "GuardNet: Graph-Attention Filtering," arXiv:2509.23037, 2025.
\bibitem{ref9} Z. Huang et al., "Attention-Aware GNN-based Input Defense (G-Guard)," arXiv:2507.07146, 2025.
\bibitem{ref10} K. Greshake et al., "Indirect Prompt Injection in LLMs," in Proc. AI  Se  CS, 2023. 
\bibitem{ref11} A. Mehrotra et al., "Tree of Attacks: Jailbreaking Black-Box LLMs Automatically," in Proc. NeurIPS, 2024.
\bibitem{ref12} P. Chao et al., "JailbreakBench: Open Robustness Benchmark," in Proc. NeurIPS Datasets, 2024.
\bibitem{ref13} Y. Feng et al., "Hypergraph Neural Networks," in Proc. AAAI, 2019.
\bibitem{ref14} S. Bai, F. Zhang, and P. H. Torr, "Hypergraph Convolution," Pattern Recognition, 2021.
\end{thebibliography}

\end{document}